\chapter{Notation and Preliminaries}\label{chap:notation}

We first explain some notation that shall be used throughout this article. 

\begin{itemize}
\item In most cases, the degree of the polynomial shall be denoted by $d$ and the number of variables shall be denoted by $n$. 
In some cases, $n$ would refer to a parameter that would determine the number of variables (though perhaps not exactly). 
\item Almost all the polynomials that we shall be studying would be multilinear. 
In multilinear polynomials, we shall identify a monomial with the \emph{set of variables} that it is a product of. 
This would allow us to abuse notation and say $x_i \in m$ to mean that $x_i$ divides the monomial $m$, or to say $m_1 \intersection m_2$ to refer to the $\gcd(m_1,m_2)$. 
We shall also use the notation $m \in f$ to mean that the polynomial $m$ has a non-zero coefficient for the monomial $m$. 
\item We shall use $[n]$ to denote $\inbrace{1,\dots, n}$ and we shall use boldface letters such as $\mathbf{x}$ to denote sets of variables. 
Further, $\mathbf{x}_{[n]}$ shall refer to a set of variables $\inbrace{x_1,\dots, x_n}$. 
However, if the number of variables is understood, we shall drop the subscript. 
\end{itemize}


\section{Polynomials of interest}

There are a few polynomials that are the usual suspects while proving lower bounds. 
The polynomials that we would be dealing with in this article are defined below. 

\subsection*{The determinant and the permanent families}

The determinant of an $n\times n$ symbolic matrix shall be denoted by $\Det_n$ and is defined as
\[
\Det_n \spaced{=} \sum_{\sigma \in S_n} \text{sign}(\sigma) x_{1,\sigma(1)}\dots x_{n,\sigma(n)}.
\]
The permanent of an $n\times n$ symbolic matrix shall be denoted by $\Perm_n$ and is defined as
\[
\Perm_n \spaced{=} \sum_{\sigma \in S_n}x_{1,\sigma(1)}\dots x_{n,\sigma(n)}.
\]

Both of these polynomials are of degree $n$ and over $n^2$ variables. 
We know that $\Det_n$ can be computed by a polynomial sized arithmetic circuit and it is widely believed that the permanent requires circuits of size $2^{\Omega(n)}$. 

\subsection*{The Nisan-Wigderson polynomial families}

\begin{definition}[Nisan-Wigderson Polynomials]\label{defn:NW-polyfamilies}
Let $n,m,d$ be arbitrary parameters with $m$ being a power of a prime, and $n,d\leq m$. 
Since $m$ is a power of a prime, let us identify the set $[m]$ with the field $\F_m$ of $m$ elements. 
Note that since $n \leq m$, we have that $[n] \subseteq \F_m$. 
The Nisan-Wigderson polynomial with parameters $n,m,d$, denoted by $\mathrm{NW}_{n,m,d}$ is defined as
\[
\NW_{n,m,d}(\vecx) \spaced{=} \sum_{\substack{p(t) \in \F_m[t]\\ \deg(p) \leq d}} x_{1,p(1)}\dots x_{n,p(n)}.
\]
That is, for every univariate polynomial $p(t) \in \F_m[t]$ of degree at most $d$, we add one monomials that encodes the `graph' of $p$ on the points $[n]$. 
This is a polynomial of degree $n$ over $mn$ variables.
\end{definition}


This monomials of this polynomial satisfy a very useful ``low pairwise-intersection'' property. 

\begin{lemma}\label{lem:NW-low-intersection}
Let $m_1$ and $m_2$ be any two distinct monomials in $\NW_{n,m,d}(\vecx)$. 
Then, there are at most $d$ variables that divide both $m_1$ and $m_2$. 
\end{lemma}
\begin{proof}
Let $m_1$ and $m_2$ correspond to univariates $p_1(t), p_2(t) \in \F_m[t]$ of degree at most $d$. 
Then if $x_{ij}$ divides $m_1$, then $p_1(i) = j$, similarly for $m_2$.
But since $p_1$ and $p_2$ are two distinct polynomials of degree at most $d$, they can agree in at most $d$ evaluations. 
Thus, there can be at most $d$ variables that divide both $m_1$ and $m_2$.
\end{proof}

For most generic choices of the parameters, the polynomial $\NW_{n,m,d}$ is believed to require circuits of exponential size to compute them. 

\subsection*{The Iterated-Matrix-Multiplication polynomial}

For parameters $n$ and $d$, the Iterated-Matrix-Multiplication polynomial, denoted by $\IMM_{n,d}$, is defined as follows
$$
\IMM_{n,d} \spaced{=} \sum_{1\leq i_1,\dots, i_d\leq n} x_{1,i_1}^{(1)}x_{i_1,i_2}^{(2)}\dots x_{i_{d-2},i_{d-1}}^{(d-1)}x_{i_{d-1},1}^{(d)}.
$$
An equivalent way of defining the polynomial as the $(1,1)$-th entry of the product of $d$ generic $n\times n$ matrices:
$$
\IMM_{n,d} \spaced{=} \inparen{\insquare{\begin{array}{ccc} x_{11}^{(1)} & \dots & x_{1n}^{(1)}\\ \vdots & \ddots & \vdots \\ x_{n1}^{(1)} & \dots & x_{nn}^{(1)}\end{array}} \cdots \insquare{\begin{array}{ccc} x_{11}^{(d)} & \dots & x_{1n}^{(d)}\\ \vdots & \ddots & \vdots \\ x_{n1}^{(d)} & \dots & x_{nn}^{(d)}\end{array}}}_{(1,1)}.
$$

It is often useful to think of this as the polynomial computed by a \emph{generic algebraic branching program} of width $n$ and depth $n$ (where the edge connecting vertex $i$ of layer $\ell$ to vertex $j$ of layer $\ell+1$ has weight $x_{ij}^{(\ell)}$). 

This is a polynomial of degree $d$ and over $n^2(d-2) + 2n$ variables\footnote{Only the first row of the first matrix, and the first column of the last matrix participates in the $(1,1)$ entry of the product}.
Further, since the polynomial corresponds to a generic algebraic branching program, $\IMM_{n,d}$ can be computed by an arithmetic circuit of size $\poly(n,d)$. 

\subsection*{The Elementary Symmetric Polynomial}

%\subsection*{Schur Polynomials (?)}

\section{Models of computation}

As mentioned earlier, the most robust model for computing polynomials are arithmetic circuits, which are formally defined as follows. 

\begin{definition}[Arithmetic circuits and formulas]\label{defn:arithmetic-circuit}
An arithmetic circuit is a directed acyclic graph with a unique sink vertex called the \emph{root}. 
The source vertices are labelled by either formal variables or field constants, and each internal node of the graph is labelled by either $+$ or $\times$. 
Nodes compute formal polynomials in the input variables  in the natural way. 
Further, edges entering a $+$ nodes also might have field constants on them to allow the $+$ to compute an $\F$-linear combination of the children (rather than just a sum). 

The polynomial computed by the circuit is defined as the polynomial computed by the root.

\noindent
If the underlying graph is a \emph{tree} instead of a general acyclic graph, the circuit is called a \emph{formula}.
\end{definition}


Another model of computation that is studied often is the model \emph{algebraic branching programs}, defined as follows. 

\begin{definition}[Algebraic Branching Program]\label{defn:ABP}
An algebraic branching program (ABP) is a layered graph with a unique source vertex (that we shall call $s$) and a unique sink vertex (that we shall call $t$). 
All edges are from layer $i$ to layer $i+1$, and each edge is labelled by a linear polynomial. 
The polynomial computed by the ABP is defined as 
\[
f\spaced{=} \sum_{\gamma: s \leadsto t} \mathrm{wt}(\gamma)
\]
where for every path $\gamma$ from $s\leadsto t$, the weight $\mathrm{wt}(\gamma)$ is defined as the product of the labels over the edges in $\gamma$. 
\end{definition}

The width of the ABP is defined as the maximum number of vertices in any layer, and the depth is defined as the length of the longest path from $s$ to $t$. 
The polynomial computed by an ABP is captured by the \emph{iterated matrix multiplication} polynomial that we shall soon see. 

It is easy to observe that any arithmetic formula of size $s$ can be simulated by an algebraic branching program of size $\poly(s)$, and any algebraic branching program of size $s$ can be simulated of an arithmetic circuit of size $\poly(s)$. 
It is a major open problem to show a separation between any of these.   
\[
    \mathrm{Formulas} \spaced{\subseteq} \mathrm{ABPs} \spaced{\subseteq} \mathrm{Circuits}
\]

\begin{openproblem}
Show a super-polynomial separation between any of the models -- formulas, ABPs and circuits. 
\end{openproblem}

\subsection{Constant depth circuits}

We shall be dealing a lot with constant depth circuits. 
Normally, the root is assumed to be a $+$ gate\footnote{The reason for this is that often the polynomial computed by the circuits would be irreducible, and hence would be silly to have a $\times$ gate as a root.} and the circuit is assumed to consist of alternating layers of $+$ and $\times$ gates.
A layer of $+$ gates are called $\Sigma$ layers, and a layer of $\times$ gates are called $\Pi$ layers. 
Thus, a depth two circuit consist of the form 
\[
f\quad = \quad \sum_{i=1}^s \prod_{j=1}^d x_{ij}
\]
is a $\Sigma\Pi$ circuit. 

\begin{fact}
Any arithmetic circuit of depth $\Delta$ and size $s$, can be simulated by an arithmetic formula of depth $\Delta$ and size $s' \leq s^{\Delta}$. 
\end{fact}

Thus for constant depth circuits (where $\Delta = O(1)$), we may assume that we are dealing with formulas without much loss of generality.


It would also be useful to keep track of the \emph{fan-in} of the gates in a certain layer (especially of multiplication gates). 
We shall use superscripts to denote this. 
For example, an $\Sigma\Pi^{[a]}\Sigma\Pi^{[b]}$ circuits computes a polynomial of the form
\[
f \spaced{=} \sum_i \prod_{j=1}^a Q_{ij}
\]
where each $Q_{ij}$ is a polynomial of degree at most $b$. \\

It would also be useful to consider special layers of multiplication gates that multiply the same polynomial several times, rather than multiplying several different polynomials together. 
Since such gates simply raise the input to a certain power, these would be called \emph{exponentiation} gates. 
A layer of exponentiation will be denoted by $\wedge$ \footnote{To say a little on the choice of notation, it was introduced in \cite{gkks13b} and the first choice was to use ${}^{\wedge}$, but looked rather ugly to write it as say $\Sigma{}^{\wedge}\Sigma$. 
Hence, $\Sigma\!\wedge\!\Sigma$  was chosen instead. } and, for example, a $\Sigma\!\wedge\!\Sigma$ circuit computes a polynomial of the form 
\[
f\spaced{=} \sum_{i=1}^s \ell_{i}^d
\]
where each $\ell_i$ is a linear polynomial. \\

\begin{exercise}
Show that any constant width ABP can be simulated by a polynomial sized formula. 
\end{exercise}

\subsection{Homogenity}\label{sec:homogenization}

Suppose we have an $n$-variate degree $d$ polynomial computed by an arithmetic circuit $C$. 
How large can the degree of intermediate computations be? 
Potentially, intermediate computations can involve very high degree terms which somehow cancel each other at the root. 
However, the following lemma of Strassen shows that we may assume without much loss of generality that arithmetic circuits never compute polynomials of degree more than the output. 

\begin{definition}[Homogeneous circuits]
A circuit $C$ is said to be \emph{homogeneous} if every gate in the circuit computes a homogeneous polynomial. 
\end{definition}

\begin{lemma}[Homogenization]\label{lem:homogenization}
    Let $f$ be an $n$-variate degree $d$ polynomial computed by a circuit $C$ of size $s$. 
    Then, for every $0\leq i \leq d$, there is a \emph{homogeneous arithmetic circuit} $C_i'$, of size at most $O(sd^2)$, that computes the degree $i$ homogeneous polynomial in $f$ (the sum of all degree $i$ monomials in $f$).  
\end{lemma}
\begin{proof}
    Assume without loss of generality that the circuit $C$ has all gates with fan-in at most $2$. 
    For every gate $g\in C$, define $(d+1)$ gates $g^{(0)},\dots, g^{(d)}$; we shall construct a new circuit $C'$ such that $g^{(i)}$ computes the degree $i$ homogeneous component of the polynomial computed at $g$. 
    If $g$ has children $h_1$ and $h_2$, then $C'$ would have the following connections depending on the type of the gate $g$:
    \begin{eqnarray*}
        \text{$g = h_1 + h_2$}\quad\implies\quad g^{(i)} &=& h_1^{(i)} \spaced{+} h_2^{(i)}\quad\text{for all $i$}\\
        \text{$g = h_1 \times h_2$}\quad\implies\quad g^{(i)} &=& \sum_{j=0}^i h_1^{(j)} h_2^{(i-j)}\quad\text{for all $i$}
    \end{eqnarray*}
    It is easy to check that the size of the circuit $C'$ is at most $O(sd^2)$, and computes all the homogeneous components of $f$. 
\end{proof}

Thus, for arithmetic circuits, we can assume without much loss of generality that we are working with a homogeneous circuit. 
The class of ABPs can also be assumed to be homogeneous without loss of generality. 
We leave this as an exercise.

\begin{exercise}
    Prove a similar homogenization lemma for algebraic branching programs. 
\end{exercise}

\begin{remark*}
    For the class of arithmetic formulas, it is not clear if we can homogeneous without loss of generality. 
    If we were to apply the above lemma to an arbitrary arithmetic formula, the resulting object is a homogeneous circuit and not a formula. 
    It is unclear if any formula can be homogenized without loss of generality. 
    The same is the case even for circuits of a fixed depth, as the above construction doubles the depth of the circuit.  
\end{remark*}


\subsection{Multilinear and Set-Multilinear Settings}

Most of the polynomials that are studied usually, like those described in \autoref{chap:notation}, are multilinear. 
A natural question is whether or not multilinear polynomials can be computed in a ``multilinear fashion''. 
This is formalized by what the model of multilinear circuits, in a way analogous to homogeneous circuits. 

\begin{definition}[Multilinear circuits]
A circuit $C$ is said to be \emph{multilinear} if every gate of the circuit computes a multilinear polynomial. 
A circuit is said to be \emph{syntactically multilinear} if for any $g = g_1 \times g_2$, there is no variable that has a path to both $g_1$ and $g_2$. 
\end{definition}

Note that syntactic multilinearity of course implies multilinearity as the definition forces all gates to compute multilinear polynomials. 
However, we could have a setting where there is a gate $g = g_1 \times g_2$ where some variable $x$ has a path to both $g_1$ and $g_2$ but it so turns out that $g_1$ is independent of $x$ due to other cancellations. 
However, for arithmetic formulas, the two notions are equivalent. 

\begin{exercise}
Given any arithmetic formula $\Phi$ that is multilinear, show that it can be converted to a formula $\Phi'$ of size $\poly(\Phi)$ that is syntactically multilinear. 

\noindent
{\bf (Imp.)} Why does the same not work for multilinear circuits? 
\end{exercise}

\begin{definition}[Set-multilinear polynomials]\label{defn:set-multilinear}
  Let $\vecx = \vecx_1 \sqcup \cdots \sqcup \vecx_d$ be a partition of variables and let $|\vecx_i| = m_i$.
A polynomial $f(\vecx)$ is said to be \emph{set-multilinear} with respect to the above partition if every monomial $m$ in $f$ satisfies $\abs{m \intersection X_i} = 1$ for all $i \in [d]$.
\end{definition}
In other words, each monomial in $f$ picks up at most one variable from each part in the partition. It is easy to see that many natural polynomials such as $\Det$ or $\Perm$ or $\NW$ are all set-multilinear for an appropriate partition of variables.

\subsection{Non-Commutative Setting}

A non-commutative polynomials over many variables, denoted by $\F\inbrace{x_1,\ldots, x_n}$, are formal polynomials over the variables wherein the variables do not commute.
Hence, a polynomial $x_1x_2 - x_2x_1$ in $\F\set{x_1,\ldots, x_n}$ is a non-zero polynomial.
They naturally can be added or multiplied but the order in which the variables are multiplied become important.
Hence,
\[
(x_1 + x_2)(x_1 + x_2) \spaced{=} x_1^2 + x_2x_1 + x_1x_2 + x_2^2 \spaced{\neq} x_1^2 + 2x_1x_2 + x_2^2
\]
Each monomial is no longer identifiable by just an exponent vector but is rather a \emph{word} on the set $\set{x_1,\ldots, x_n}$. 

In this space, we can continue to talk about arithmetic circuits or algebraic branching programs where we always keep track of the order of variables multiplied.
In arithmetic circuits or formulas, every $\times$ gate has labelled left and right children.
In an algebraic branching program, the weight of a path from source to sink is the product of the edge weights \emph{in the order from left to right}.

\subsection{Monotone Setting}

A monotone circuit over $\RR$ is a circuit having $+, \times$ gates in which all the field constants are {\em non-negative} real numbers. 
Such a circuit can compute any polynomial $f$ over $\RR$ all of whose coefficients are nonnegative real numbers, such as for example the permanent. 
It is then natural to ask whether there are small monotone circuits over $\RR$ computing the permanent. 
Jerrum and Snir \cite{js82} obtained an exponential lower bound on the size of monotone circuits over $\RR$ computing the permanent. 
Note that this definition of monotone circuits is valid only over $\RR$ (actually more generally over ordered fields but not over say finite fields) and such circuits can only compute polynomials with non-negative coefficients. 
Here we will present Jerrum and Snir's argument in a slightly more generalized form such that the circuit model makes sense over any field $\FF$ and is complete, i.e.  can compute any polynomial over $\FF$.\footnote{This generalization was told to me by Neeraj Kayal} Let us first explain the motivation behind the generalized circuit model that we present here. 
Observe that in any monotone circuit over $\RR$, there is no cancellation as there are no negative coefficients. 
Formally, for a node $v$ in our circuits let us denote by $f_{v}$ the polynomial computed at that node. 
For a polynomial $f$ let us denote by $\Mon(f)$ the set of monomials having a nonzero coefficient in the polynomial $f$.
\begin{enumerate}
\item If $w = u + v$ then 
  $$ \Mon(f_{w}) = \Mon(f_{u}) \cup \Mon(f_{v}). $$
  
\item If $w = u \times v$ then
  $$ \Mon(f_{w}) = \Mon(f_{u}) \cdot \Mon(f_{v}) \eqdef \setdef{m_1\cdot m_2}{m_1 \in \Mon(f_{u}), m_2\in \Mon(f_{v})}. $$
			
\end{enumerate}
This means that for any node $v$ in a monotone circuit over $\RR$ one can determine $\Mon(f_{v})$ in a very syntactic manner starting from the leaf nodes. 
Let us make precise this syntactic computation that we have in mind.

\begin{definition}[Formal Monomials.]
  Let $\Phi$ be an arithmetic circuit. 
The \emph{formal monomials} at any node $v\in \Phi$, which shall be denoted by $\FM(v)$, shall be inductively defined as follows:
  \begin{quote}
    If $v$ is a leaf labelled by a variable $x_i$, then $\FM(v) = \inbrace{x_i}$. 
If it is labelled by a constant, then $\FM(v) = \inbrace{1}$.
    
    If $v = v_1 + v_2$, then $\FM(v) = \FM(v_1) \union \FM(v_2)$. 
    
    If $v = v_1 \times v_2$, then 
    \begin{eqnarray*}
      \FM(v) &=& \FM(v_1)\cdot \FM(v_2)\\
      &\eqdef& \setdef{m_1\cdot m_2}{m_1 \in \FM(v_1), m_2\in \FM(v_2)}.
    \end{eqnarray*}
  \end{quote}
\end{definition}

\noindent Note that for any node $v$ in any circuit we have $\Mon(f_{v}) \subseteq \FM(v)$ but in a monotone circuit over $\RR$ this containment is in fact an equality at every node. 
This motivates our definition of a slightly more general notion of a monotone circuit as follows.

\begin{definition}[Monotone circuits]\label{defn:monotone-circuit}
  A circuit $C$ is said to be \emph{syntactically monotone} (simply monotone for short) if $\Mon(f_{v}) = \FM(v)$ for every node $v$ in $C$.
\end{definition}
